%\clearpage
\section*{Impact Statements}
We note that the construction of \simulator{} does not appear to increase the probability of harmful behavior, because the construction's primary objective is to implicitly tune an internal model (\textsection{\ref{sec:design}}).
Such tuning has been possible for a long time and is not made more expressive by \simulator{}.

Our findings suggest that existing transformer-based language models can plausibly possess the ability to learn and adapt to context by internally fine-tuning a complex model \emph{even during inference}. 
Consequently, although users are unable to directly modify deployed models, these models may still undergo dynamic updates while processing a context left-to-right, resulting in previously unseen behavior by the time the model reaches the end of the context.
This has significant implications for the field of model alignment.
It is challenging to impose restrictions on a model that can perform such dynamics updates internally, so malicious content can influence the output of deployed models.


Alternatively, we recognize the potential benefits of pre-training constructed models that integrate explicit fine-tuning mechanisms. By embedding the functionalities typically achieved through explicit fine-tuning, such as detecting malicious content and intent within the models themselves, the need for external modules can be mitigated. Pre-training the constructed model may offer a self-contained solution for ensuring safe and responsible language processing without relying on external dependencies.


\section*{Acknowledgements}
The authors acknowledge funding from NSF, ONR, Simons Foundation, and DARPA.
We thank Danqi Chen, Jason Lee, Zhiyuan Li, Kaifeng Lyu, Simran Kaur, Tianyu Gao, and Colin Wang for their suggestions and helpful discussions at different stages of our work.